\documentclass[11pt,a4paper]{article}
\usepackage[utf8]{inputenc}
\usepackage[margin=1in]{geometry}
\usepackage{amsmath,amssymb,amsfonts}
\usepackage{booktabs}
\usepackage{xcolor}
\usepackage{hyperref}
\usepackage{microtype}
\usepackage{array}

\definecolor{primary}{RGB}{31, 78, 121}
\definecolor{secondary}{RGB}{84, 130, 53}

\hypersetup{
    colorlinks=true,
    linkcolor=primary,
    urlcolor=primary
}

\title{\textbf{\Large Axomeme Architecture Specification}\\[0.5em]
\large A Self-Contained Mathematical and Algorithmic Specification for Neural Detection of Natural Selection in Molecular Sequences}
\author{\textbf{Axomeme Deep Learning Engineering Team}}
\date{August 10, 2026}

\begin{document}
\maketitle

\begin{abstract}
This document provides the authoritative mathematical and algorithmic specification of the \textbf{Axomeme} neural network architecture. Axomeme is an end-to-end deep learning framework designed to predict site-level natural selection (Likelihood Ratio Test statistics) directly from multi-species codon sequence alignments and phylogenetic tree matrices. The architecture combines a 4-layer dual-axial transformer engine with an evolutionarily grounded \textbf{Markov Exponential Decay Kernel} ($P(t) = e^{-\lambda t}$), 4D metric tree Multidimensional Scaling (MDS), block-diagonal track disentanglement between synonymous and non-synonymous features, a \textbf{12-Expert Mixture-of-Pooling Streams (MoPS)} gating network, and a rank-consistent CORAL ordinal prediction head.
\end{abstract}

\tableofcontents
\newpage

\section{Introduction \& Problem Formulation}

\subsection{The Statistical Framework of Natural Selection}
In evolutionary biology, protein-coding gene sequences evolve under varying selective pressures:
\begin{enumerate}
    \item \textbf{Synonymous Substitutions ($\alpha$ or $dS$):} Nucleotide changes that do not alter the translated amino acid.
    \item \textbf{Non-Synonymous Substitutions ($\beta$ or $dN$):} Nucleotide changes that alter the translated amino acid.
\end{enumerate}

Site-level selection is formulated as a hypothesis test:
\begin{align}
H_0 \text{ (Null Hypothesis)}&: \beta \le \alpha \quad (\text{Neutral or Purifying Selection}) \\
H_1 \text{ (Alternative Hypothesis)}&: \beta^+ > \alpha \quad (\text{Episodic Positive Selection on a subset of branches})
\end{align}

The Likelihood Ratio Test (LRT) statistic is defined as:
\begin{equation}
\text{LRT} = 2 \cdot \Big(\log \mathcal{L}(H_1) - \log \mathcal{L}(H_0)\Big)
\end{equation}

Axomeme models the mapping from raw multi-species sequence alignments and tree matrices directly to LRT statistics, replacing numerical likelihood optimization with a direct neural operator.

---

\section{Input Representation \& Encoding}

Consider an alignment matrix $X$ containing $N$ species sequences across $S$ codon sites.

\subsection{Discrete Sequence Tokenization}
The alignment is tokenized into two discrete integer matrices:
\begin{enumerate}
    \item \textbf{Codon Tokens ($X_{\text{codon}} \in \{0, \dots, 65\}^{N \times S}$):} 64 sense codons + 2 termination/gap tokens.
    \item \textbf{Amino Acid Tokens ($X_{\text{aa}} \in \{0, \dots, 22\}^{N \times S}$):} 20 standard amino acids (indices 0--19) + 1 stop codon (20) + 1 gap token (21) + 1 unknown token (22).
\end{enumerate}

\subsection{Phylogenetic Distance \& Coordinates}
\begin{enumerate}
    \item \textbf{Pairwise Patristic Distance Matrix ($\mathbf{T}_{\text{phylo}} \in \mathbb{R}^{N \times N}$):} Patristic branch distances are computed from the unrooted phylogenetic tree:
    \begin{equation}
    \tau_{ij} = \sum_{b \in \text{Path}(i, j)} \text{Length}(b)
    \end{equation}
    \item \textbf{Pairwise Jukes-Cantor Distance Fallback:} If branch lengths are absent, pairwise distances are calculated directly from sequence identity:
    \begin{equation}
    d(s_i, s_j) = -\frac{3}{4} \ln\left(1 - \frac{4}{3} p_{ij}\right)
    \end{equation}
    \item \textbf{4D Multidimensional Scaling (MDS) Coordinates ($M \in \mathbb{R}^{N \times 4}$):} The distance matrix $\mathbf{T}_{\text{phylo}}$ is projected into 4-dimensional Euclidean spatial coordinates via classical metric MDS.
\end{enumerate}

---

\section{Block-Diagonal Track Disentanglement}

To isolate non-synonymous amino acid selection signals from synonymous nucleotide noise, Axomeme enforces \textbf{Track Disentanglement} throughout the embedding and transformation layers.

\subsection{Vector Partitioning}
Embedding vectors $\mathbf{x} \in \mathbb{R}^{256}$ are partitioned into two independent 128-dimensional sub-vectors:
\begin{equation}
\mathbf{x} = \begin{bmatrix} \mathbf{x}_{\text{codon}} \\ \mathbf{x}_{\text{aa}} \end{bmatrix} \in \mathbb{R}^{256}, \quad \mathbf{x}_{\text{codon}} \in \mathbb{R}^{128}, \quad \mathbf{x}_{\text{aa}} \in \mathbb{R}^{128}
\end{equation}

\subsection{The BlockLinear Operator}
All linear projection layers use a block-diagonal weight matrix \texttt{BlockLinear}:
\begin{equation}
\mathbf{W}_{\text{block}} = \begin{pmatrix} \mathbf{W}_{\text{codon}} & \mathbf{0} \\ \mathbf{0} & \mathbf{W}_{\text{aa}} \end{pmatrix} \in \mathbb{R}^{256 \times 256}
\end{equation}

\begin{equation}
\mathbf{W}_{\text{block}} \mathbf{x} = \begin{bmatrix} \mathbf{W}_{\text{codon}} \mathbf{x}_{\text{codon}} \\ \mathbf{W}_{\text{aa}} \mathbf{x}_{\text{aa}} \end{bmatrix}
\end{equation}

---

\section{The Dual-Axial Attention Engine \& Markov Exponential Decay Kernel}

\subsection{Species Row Attention \& Markov Distance Kernel}
For each attention head $h \in \{1, \dots, 8\}$:
\begin{enumerate}
    \item \textbf{Raw Query-Key Similarity Score:}
    \begin{equation}
    S_{ij}^{(h)} = \frac{(\mathbf{q}_i^{(h)})^\top \mathbf{k}_j^{(h)}}{\sqrt{32}}
    \end{equation}
    \item \textbf{Markov Continuous-Time Decay Kernel ($\tilde{S}_{ij}^{(h)} \in \mathbb{R}$):}
    Under continuous-time Markov substitution models ($P(t) = e^{\mathbf{Q}t}$), evolutionary influence between species separated by patristic distance $\tau_{ij}$ decays exponentially: $K(\tau_{ij}) \propto e^{-\lambda_h \cdot \tau_{ij}}$. In attention score logit space, this corresponds to subtracting a learnable positive distance penalty:
    \begin{equation}
    \tilde{S}_{ij}^{(h)} = \frac{(\mathbf{q}_i^{(h)})^\top \mathbf{k}_j^{(h)}}{\sqrt{32}} - \lambda_h \cdot \tau_{ij}, \quad \text{where } \lambda_h = \text{Softplus}(w_{\text{phylo}, h}) > 0
    \end{equation}
    \item \textbf{Softmax Normalization:}
    \begin{equation}
    A_{ij}^{(h)} = \frac{\exp(\tilde{S}_{ij}^{(h)})}{\sum_{m=1}^N \exp(\tilde{S}_{im}^{(h)})}
    \end{equation}
\end{enumerate}

---

\section{Mixture-of-Pooling Streams (MoPS) Gating Architecture}

Following 4 layers of phylogenetic species attention, each species $i$ at site $s$ has a 256-dimensional embedding vector $\mathbf{h}_i \in \mathbb{R}^{256}$. Axomeme pools species representations across $K=12$ specialist experts:

\begin{table}[h]
\centering
\caption{The 12 MoPS Specialist Pooling Experts}
\vspace{0.5em}
\begin{tabular}{l c l p{5.5cm}}
\toprule
\textbf{Expert Stream} & \textbf{Dims} & \textbf{Mathematical Definition} & \textbf{Primary Biological Function} \\
\midrule
1. \texttt{mean\_pooled} & 256 & $\boldsymbol{\mu} = \frac{1}{N} \sum_{i=1}^N \mathbf{h}_i$ & Equilibrium Consensus Background \\
2. \texttt{max\_pooled} & 256 & $\mathbf{m} = \max_{i} \mathbf{h}_i$ & Deep Mutant Outlier Burst \\
3. \texttt{min\_pooled} & 256 & $\mathbf{n} = \min_{i} \mathbf{h}_i$ & Purifying Constraint Floor \\
4. \texttt{cat\_pooled} & 256 & $\text{MLP}_{\text{fusion}}\Big( \Big[ \text{MLP}_{\text{stream}}\Big( \mathbf{e}_a \,||\, p_a \Big) \Big]_{a=0}^{19} \Big)$ & Clean 20-Class Allele Matrix \\
5. \texttt{diff\_pooled} & 256 & $\text{ReLU}(\mathbf{m} - \boldsymbol{\mu}) \in \mathbb{R}^{256}$ & Full 256D Codon + AA Disentangled Range \\
6. \texttt{var\_pooled} & 256 & $\frac{1}{N} \sum_{i=1}^N (\mathbf{h}_i - \boldsymbol{\mu})^2 \in \mathbb{R}^{256}$ & 2nd Moment (Feature Variance) \\
7. \texttt{skew\_pooled} & 256 & $\frac{1}{N} \sum_{i=1}^N (\mathbf{h}_i - \boldsymbol{\mu})^3 \in \mathbb{R}^{256}$ & 3rd Moment (Asymmetric Lineage Acceleration) \\
8. \texttt{kurt\_pooled} & 256 & $\frac{1}{N} \sum_{i=1}^N (\mathbf{h}_i - \boldsymbol{\mu})^4 \in \mathbb{R}^{256}$ & 4th Moment (Rare Outlier Tailness) \\
9. \texttt{lse\_pooled} & 256 & $\log \big( \frac{1}{N} \sum_{i=1}^N \exp(\mathbf{h}_i) \big)$ & Smooth Log-Sum-Exp Max/Mean Interpolation \\
10. \texttt{phylo\_mean\_pooled} & 256 & $\frac{\sum_{i=1}^N w_i \mathbf{h}_i}{\sum_{i=1}^N w_i}, \quad w_i = \frac{1}{d_i + 1.0}$ & Tree-Weighted Consensus \\
11. \texttt{dnds\_contrast} & 256 & $\frac{1}{N} \sum_{i=1}^N \left( \frac{\|\mathbf{h}_{i, \text{aa}}\|}{\|\mathbf{h}_{i, \text{codon}}\| + \epsilon} \right) \mathbf{h}_i$ & Non-synonymous vs Synonymous Channel Norm Ratio \\
12. \texttt{phylo\_var\_pooled} & 256 & $\frac{\sum_{i=1}^N w_i (\mathbf{h}_i - \boldsymbol{\mu}_{\text{phylo}})^2}{\sum_{i=1}^N w_i}$ & Tree-Weighted Lineage Variance \\
\bottomrule
\end{tabular}
\end{table}

\subsection{Dynamic Neural Routing}
The gating network computes Softmax routing weights over the 12 expert streams based on site consensus:
\begin{equation}
\mathbf{g} = \text{Softmax}\left( \mathbf{W}_g \boldsymbol{\mu} + \mathbf{b}_g \right) \in [0, 1]^{12}, \quad \sum_{k=1}^{12} g_k = 1.0
\end{equation}

\begin{equation}
\mathbf{h}_{\text{pooled}} = \text{stream\_fusion}\left( \sum_{k=1}^{12} g_k \cdot \mathbf{S}_k \right) \in \mathbb{R}^{256}
\end{equation}

---

\section{Rank-Consistent CORAL Ordinal Head}

The pooled site vector $\mathbf{h}_{\text{pooled}} \in \mathbb{R}^{256}$ is passed to the Consistent Rank Logits (CORAL) ordinal prediction head.

\subsection{Non-Linear Projection \& Strict Logit Monotonicity}
To prevent 1D feature collapse, features are projected through a 2-layer MLP projection $\text{proj}(\mathbf{h}_{\text{pooled}})$ before adding monotonic threshold cutoffs $b_k$:
\begin{equation}
\text{proj}(\mathbf{h}_{\text{pooled}}) = \mathbf{W}_2 \cdot \text{ReLU}(\mathbf{W}_1 \mathbf{h}_{\text{pooled}}), \quad \mathbf{W}_1 \in \mathbb{R}^{64 \times 256}, \quad \mathbf{W}_2 \in \mathbb{R}^{1 \times 64}
\end{equation}

Threshold cutoffs $b_k$ enforce strict probability monotonicity ($P(y > b_0) \ge P(y > b_1) \ge \dots$):
\begin{equation}
b_0 = \text{Learnable Parameter}, \quad b_k = b_0 - \sum_{j=1}^k \text{Softplus}(\theta_j)
\end{equation}

The cumulative logits for ordinal threshold $k \in \{0, \dots, 11\}$ are:
\begin{equation}
z_k(\mathbf{h}_{\text{pooled}}) = \text{proj}(\mathbf{h}_{\text{pooled}}) + b_k
\end{equation}

\subsection{Expected Value Decoding}
Continuous LRT predictions are decoded directly via Smooth Expected Value Interpolation over the empirical MEME $\chi^2$ mixture threshold midpoints:
\begin{equation}
\hat{\text{LRT}} = \sum_{m=0}^{12} P(\text{Bin } m) \cdot \text{Midpoint}_m
\end{equation}
where $\text{Midpoint}_m$ represents the raw LRT midpoint of bin $m$, providing continuous, un-biased selection strength estimates.

\end{document}
